% Source: https://artofproblemsolving.com/community/c3h2490513_2017_amc_10b_problemsproblem_17
% Stable identity: aops-topic-2490513-q01
A strictly monotone number is determined by its set of digits, since the order is forced.
\begin{itemize}
\item \textbf{Increasing (including one-digit numbers).} The digit $0$ cannot appear: it would have to come first. Every nonempty subset of $\{1,\dots,9\}$ gives one number: $2^9-1=511$ numbers.
\item \textbf{Decreasing with at least two digits.} Any subset of $\{0,1,\dots,9\}$ with at least two elements gives one number; its first digit is its largest, which is nonzero. There are $2^{10}-1-10=1013$ such subsets.
\end{itemize}
The two lists do not overlap, since numbers with two or more digits cannot be both increasing and decreasing, and the one-digit numbers were counted only in the first list. The total is $511+1013=\boxed{\textbf{(B)}\ 1524}$.
